\documentclass[12pt,a4paper]{article}
\usepackage[T1]{fontenc}
\usepackage[margin=1in]{geometry}
\usepackage{hyperref}
\usepackage{booktabs}
\usepackage{enumitem}
\usepackage{fancyhdr}
\usepackage{longtable}
\usepackage{xcolor}

\pagestyle{fancy}
\fancyhf{}
\rhead{\thepage}
\lhead{OghmaNotes -- SRS v2.0}

\title{Software Requirements Specification \\
\Large OghmaNotes: RAG-Powered Learning Platform \\
\small Version 2.0 (Updated 2025-03-06)}
\author{Samuel Regan, Semyon Fox, Shreyansh Singh}
\date{\today}

\hypersetup{
    pdftitle={OghmaNotes SRS v2.0},
    pdfauthor={Samuel Regan, Semyon Fox, Shreyansh Singh}
}

\begin{document}

\maketitle

\tableofcontents
\newpage

%==============================================================================
\section{Introduction}
%==============================================================================

\subsection{Purpose}

This Software Requirements Specification (SRS) defines the functional and non-functional requirements for \textbf{OghmaNotes}, an AI-powered study platform combining portable Markdown notes with Retrieval-Augmented Generation (RAG) to provide intelligent learning from course materials.

This is the \emph{third iteration} of the SRS, reflecting lessons learned as the tech stack evolved to better suit the team's expertise and project goals. Rather than adhering rigidly to the original spec, we've embraced an Agile approach: document what we know, adapt when reality dictates, and deliver value within the semester deadline.

\textit{Intended Audience:} Development team (3 engineers), academic stakeholders, future maintainers, and QA.

\subsection{Document Status}

\begin{itemize}
    \item \textbf{Version:} 2.0 (Living Document -- Updated as tech stack evolves)
    \item \textbf{Last Updated:} 2025-03-06
    \item \textbf{Current Phase:} Core infrastructure complete; RAG pipeline implementation in progress
    \item \textbf{Revision History:} See Appendix C
\end{itemize}

\subsection{Scope \& Objectives}

OghmaNotes is a \textbf{web-based learning platform} allowing students to:
\begin{itemize}
    \item ✅ Create and manage Markdown notes with offline support (PWA)
    \item ✅ Upload and annotate PDFs
    \item 🔄 Ask questions via RAG chat with citations (in progress)
    \item 🔄 Index PDFs for semantic retrieval (in progress)
    \item ⏳ Generate adaptive quizzes
    \item ⏳ Use spaced-repetition flashcards
    \item ⏳ Track progress with analytics
    \item ⏳ Integrate with Canvas LMS for deadlines
\end{itemize}

\subsubsection{MVP (Must Have)}
\begin{itemize}
    \item Web-based Markdown editor with offline support (PWA)
    \item PDF ingestion, annotation, and semantic search
    \item User authentication (email/password with JWT)
    \item Note organization (folder hierarchy)
    \item RAG chat with cited answers
    \item Multi-mode search (keyword + semantic)
    \item Responsive design (desktop \& tablet)
\end{itemize}

\subsubsection{Out of Scope (MVP)}
\begin{itemize}
    \item Light mode / full WCAG AA compliance (dark mode only)
    \item Real-time multi-user editing
    \item Native mobile apps (PWA only)
    \item Instructor dashboards
    \item Video/audio transcription
\end{itemize}

\subsubsection{Key Objectives}
\begin{enumerate}
    \item Deliver functional MVP by semester end (6 weeks total)
    \item Enable 3-person parallel development via clear API boundaries
    \item Deploy automatically from \texttt{prod} branch via AWS Amplify
    \item Keep development within free tier + modest budget
\end{enumerate}

%==============================================================================
\section{Tech Stack (Actual Implementation)}
%==============================================================================

This section documents \emph{what we're actually using}, not what the original spec imagined. Stack choices were driven by team expertise, real-world viability, and pragmatism.

\subsection{Frontend}
\begin{itemize}
    \item \textbf{Framework:} Next.js 16 (App Router, SSR/SSG)
    \item \textbf{Language:} TypeScript + JSX (type safety)
    \item \textbf{UI:} React 19 + Tailwind CSS 4
    \item \textbf{Editor:} Lexical (rich text, better UX than Markdown-only)
    \item \textbf{PDF Viewer:} PDF.js + react-pdf (client-side)
    \item \textbf{State:} Zustand (lightweight)
    \item \textbf{Theme:} Dark mode only (MVP)
\end{itemize}

\subsection{Backend (Next.js API Routes)}
\begin{itemize}
    \item \textbf{Framework:} Next.js 16 (no separate backend service)
    \item \textbf{Auth:} bcryptjs (hashing) + JWT (tokens)
    \item \textbf{Email:} Nodemailer + AWS SES
    \item \textbf{Database:} PostgreSQL 12+ (postgres.js driver)
    \item \textbf{Storage:} AWS S3 (file uploads, PDFs, exports)
    \item \textbf{Middleware:} Custom CORS, session validation
\end{itemize}

\subsection{Database \& Search}
\begin{itemize}
    \item \textbf{Primary:} PostgreSQL 12+ (RDS or local)
    \item \textbf{Vector Search:} pgvector extension (1536-dim embeddings)
    \item \textbf{Full-Text Search:} PostgreSQL FULLTEXT (not yet implemented)
    \item \textbf{Migrations:} Manual SQL files (in \texttt{database/} folder)
\end{itemize}

\subsection{LLM \& AI (User-Managed Keys)}

\textbf{Key Innovation:} OghmaNotes does \textit{not} store LLM API keys. Users provide their own.

\begin{itemize}
    \item \textbf{Embedding Models:} User's choice (OpenAI, Hugging Face, Qwen, etc.)
    \item \textbf{LLM Providers:} OpenAI, Anthropic, LiteLLM gateway, custom servers
    \item \textbf{Key Management:} Users manage keys in browser (secure localStorage)
    \item \textbf{Backend Role:} Coordinates search; does \textit{not} call LLMs
    \item \textbf{Cost Model:} Users pay their provider directly
    \item \textbf{Team Testing:} Internal Qwen embedding model on dev server (free)
\end{itemize}

\textit{See Appendix B: LLM Strategy for full architecture details.}

\subsection{Storage \& External Services}
\begin{itemize}
    \item \textbf{File Storage:} AWS S3 (PDFs, Markdown, exports)
    \item \textbf{Email:} AWS SES (verification, password reset)
    \item \textbf{Caching:} AWS ElastiCache (future, for job queues)
    \item \textbf{Deployment:} AWS Amplify (auto-deploy from \texttt{prod} branch)
    \item \textbf{PDF Processing:} pdf-parse (text extraction)
    \item \textbf{OCR:} TBD (Google Vision or AWS Textract, deferred to Phase 2)
\end{itemize}

\subsection{Local Development}
\begin{itemize}
    \item \textbf{Docker:} \texttt{docker-compose up} (PostgreSQL + app)
    \item \textbf{Database:} PostgreSQL container (pgvector pre-installed)
    \item \textbf{File Storage:} AWS S3 (configure credentials in \texttt{.env.local})
    \item \textbf{Email:} Nodemailer (local SMTP or AWS SES test)
\end{itemize}

%==============================================================================
\section{Deviations from Original SRS}
%==============================================================================

The original SRS (Version 1.0) specified MariaDB and microservices. We've adapted based on team strengths and project reality. This is \textbf{not a failure to spec}---it's Agile: respond to change while staying aligned with goals.

\subsection{Database: MariaDB \textrightarrow\ PostgreSQL + pgvector}

\begin{tabular}{@{}ll@{}}
\toprule
\textbf{Original} & MariaDB (native vector columns) \\
\textbf{Actual} & PostgreSQL 12+ + pgvector extension \\
\textbf{Why Change} & Team expertise in PostgreSQL; pgvector is viable and well-maintained \\
\textbf{Impact} & Zero---both support vector search; queries slightly different syntax \\
\bottomrule
\end{tabular}

\subsection{Architecture: Microservices \textrightarrow\ Next.js Monolith}

\begin{tabular}{@{}ll@{}}
\toprule
\textbf{Original} & 3 services (frontend, API, RAG pipeline) + workers \\
\textbf{Actual} & Single Next.js app (frontend + API routes) \\
\textbf{Why Change} & Faster iteration for MVP; cleaner deployment \\
\textbf{Trade-off} & Less independent scaling; acceptable for semester project \\
\bottomrule
\end{tabular}

\subsection{Background Jobs: BullMQ + Redis \textrightarrow\ TBD}

\begin{tabular}{@{}ll@{}}
\toprule
\textbf{Original} & BullMQ + Redis (required for async) \\
\textbf{Actual} & Not yet implemented; deferred to Phase 2 \\
\textbf{Why Defer} & Long-running tasks (PDF indexing) can start as synchronous; add async later \\
\textbf{Impact} & May slow down large PDF indexing; acceptable trade-off \\
\bottomrule
\end{tabular}

\subsection{LLM Key Management: Backend-Managed \textrightarrow\ User-Managed}

\begin{tabular}{@{}ll@{}}
\toprule
\textbf{Original} & Backend stores and manages API keys \\
\textbf{Actual} & Users provide own keys (frontend only); backend never sees them \\
\textbf{Why Better} & Eliminates secret management; supports any provider; better privacy \\
\textbf{Impact} & Users pay their own providers (transparent cost model) \\
\bottomrule
\end{tabular}

%==============================================================================
\section{Functional Requirements}
%==============================================================================

Requirements organized by priority tier and status.

\subsection{Tier Definitions}

\begin{table}[h]
\centering
\begin{tabular}{@{}lll@{}}
\toprule
\textbf{Tier} & \textbf{Definition} & \textbf{MVP?} \\
\midrule
\textbf{Must Have} & Core functionality for MVP & Yes \\
\textbf{Should Have} & Important features if time permits & Maybe \\
\textbf{Nice to Have} & Polish and future work & No \\
\bottomrule
\end{tabular}
\end{table}

\subsection{Authentication \& User Management (Must Have)}

\textbf{Status:} ✅ Complete

\begin{itemize}
    \item ✅ Register with email/password
    \item ✅ Send verification email (AWS SES)
    \item ✅ Require email verification (soft enforcement)
    \item ✅ Issue JWT tokens (24h expiry)
    \item ✅ Login with clear error messages
    \item ✅ Password reset with email link (1h token)
    \item ✅ Logout / session invalidation
    \item ✅ Bcrypt hashing (10+ salt rounds)
\end{itemize}

\subsection{Note Management (Must Have)}

\textbf{Status:} ✅ Core complete; enhancements pending

\begin{itemize}
    \item ✅ Create, edit notes in Markdown
    \item ✅ Auto-save to IndexedDB (every 5s)
    \item ✅ Folder hierarchy organization
    \item ✅ Sync to S3 when online
    \item ⏳ Bidirectional links (wiki-style)
    \item ⏳ Tag-based organization
    \item ⏳ Full-text search (keyword)
\end{itemize}

\subsection{PDF \& Document Management (Must Have)}

\textbf{Status:} ⚠️ Partial (upload + annotations ✅; indexing in progress)

\begin{itemize}
    \item ✅ Upload PDFs up to 50MB
    \item ✅ PDF annotations (highlighting, notes)
    \item ⏳ Text extraction (pdf-parse)
    \item ❌ Semantic chunking (500 tokens, 50-100 overlap)
    \item ❌ Generate embeddings (OpenAI)
    \item ❌ Async job processing
    \item ⏳ OCR for scanned PDFs (future)
\end{itemize}

\subsection{RAG Chat with Citations (Must Have --- CRITICAL)}

\textbf{Status:} 🔄 In progress (RAG pipeline commit expected 2025-03-06)

\begin{itemize}
    \item ❌ Chat interface for questions
    \item ❌ Hybrid search (keyword + semantic)
    \item ❌ Retrieve relevant chunks
    \item ❌ Send context to LLM with user's key
    \item ❌ Generate answers with citations
    \item ❌ Stream responses token-by-token
    \item ❌ Rating feedback (thumbs up/down)
\end{itemize}

\textbf{Blocker for MVP:} This is the core value proposition. Highest priority for Phase 2.

\subsection{Quiz Generation (Should Have)}

\textbf{Status:} ❌ Not started; Phase 3

\begin{itemize}
    \item Generate 5-10 questions from notes/PDFs
    \item Multiple-choice, short-answer, T/F types
    \item Optional timed mode
    \item Display score + explanations
    \item Track weak topics
\end{itemize}

\subsection{Spaced Repetition Flashcards (Should Have)}

\textbf{Status:} ❌ Not started; Phase 3

\begin{itemize}
    \item Create decks from notes
    \item Display daily review queue
    \item Rate recall (1-5 scale)
    \item Update SM-2 intervals
    \item Persist state across devices
\end{itemize}

\subsection{Offline Support (Must Have)}

\textbf{Status:} ⚠️ Partial (PWA framework ready; sync queue TBD)

\begin{itemize}
    \item ✅ Installable as PWA
    \item ⚠️ Full note editing offline
    \item ⏳ Queue changes for sync
    \item ⏳ Auto-sync on reconnect (<30s)
    \item ⏳ Conflict resolution (last-write-wins)
\end{itemize}

\subsection{LMS Integration (Should Have)}

\textbf{Status:} ❌ Not started; Phase 3

\begin{itemize}
    \item OAuth connection to Canvas
    \item Import assignments + due dates
    \item Display in calendar
    \item Auto-schedule study sessions
    \item Daily sync + manual refresh
\end{itemize}

\subsection{Search (Must Have) --- Fuzzy + Semantic}

\textbf{Status:} 🔄 Architecture planned (feature/search branch); implementation beginning Phase 1

\textbf{Design:} Multi-layered search combining fuzzy full-text search + semantic vector search with tree sort capabilities. Keyboard-first UI (Cmd+K overlay).

\begin{itemize}
    \item ✅ Planned: Fuzzy full-text search (PostgreSQL FTS with typo tolerance)
    \item ✅ Planned: Semantic search via vector embeddings (pgvector ivfflat index)
    \item ✅ Planned: Search by filename + extracted text
    \item ✅ Planned: Relevance scoring for both search types
    \item ✅ Planned: Cmd+K overlay UI (keyboard-first, power-user friendly)
    \item ✅ Planned: Pagination (offset-limit, 20 results per page)
    \item ✅ Planned: Tree sorting (alphabetical, recent)
    \item ⏳ Future: Date range filtering, saved filters (MVP 2)
    \item ⏳ Future: Full-text highlighting in results (MVP 3)
\end{itemize}

\textbf{Database Schema:} New columns on \texttt{app.notes}: \texttt{extracted\_text}, \texttt{embedding} (vector 1536), \texttt{relative\_path}, \texttt{search\_vector}. Indexes: GIN (full-text), ivfflat (semantic), compound (user+date).

\textbf{See:} \texttt{SEARCH\_ARCHITECTURE\_PLAN.md} for full implementation details, API endpoints, and scalability analysis.

\subsection{Other Features}

\begin{table}[h]
\centering
\begin{tabular}{@{}lll@{}}
\toprule
\textbf{Feature} & \textbf{Tier} & \textbf{Status} \\
\midrule
Calendar/Timetable & Should & ❌ Phase 3 \\
Analytics Dashboard & Should & ❌ Phase 3 \\
Vault Export (ZIP) & Should & ⏳ Phase 2 \\
Command Palette (Cmd+K Search) & Must & ✅ Phase 1 (part of search) \\
Study Group Sharing & Nice & ❌ Deferred \\
i18n Support & Nice & ⏳ Framework ready \\
Note Export (PDF/CSV) & Nice & ⏳ Phase 3 \\
\bottomrule
\end{tabular}
\end{table}

%==============================================================================
\section{Non-Functional Requirements}
%==============================================================================

\subsection{Performance}

\begin{longtable}{@{}lll@{}}
\toprule
\textbf{Requirement} & \textbf{Target} & \textbf{Status} \\
\midrule
\endhead
RAG end-to-end response & $\leq 3$s & ⏳ TBD (Phase 2) \\
PDF indexing (50MB) & $\leq 2$min & ⏳ TBD \\
Keyword search & $\leq 5$ms & ⏳ TBD \\
Semantic search & $\leq 50$ms & ⏳ TBD \\
Hybrid search & $\leq 100$ms & ⏳ TBD \\
Note sync (500KB) & $\leq 2$s & ✅ ~500ms observed \\
Frontend cold load & $\leq 4$s & ✅ ~2.5s observed \\
Frontend cached load & $\leq 2$s & ✅ ~1s observed \\
\bottomrule
\end{longtable}

\subsection{Security}

\begin{itemize}
    \item ✅ JWT required on protected endpoints
    \item ✅ Bcrypt hashing (10+ rounds)
    \item ✅ TLS 1.2+ (HTTPS only)
    \item ✅ CORS policy enforced
    \item ⏳ Rate limiting (deferred, not critical for MVP)
    \item ✅ Credentials in environment variables only
    \item ✅ S3 presigned URLs (1h expiry)
    \item ✅ User-scoped database queries
    \item ✅ Password reset tokens (1h, single-use)
    \item ✅ Email verification tokens (24h)
\end{itemize}

\subsection{Reliability \& Availability}

\begin{itemize}
    \item ✅ 99\% availability (AWS Amplify SLA)
    \item ✅ Automated database backups (RDS)
    \item ⏳ Job retry logic (when jobs implemented)
    \item ⚠️ Graceful degradation (basic; can improve)
    \item ✅ Session persistence across restarts
    \item ⏳ PWA reconnection (<30s)
\end{itemize}

\subsection{Usability}

\begin{itemize}
    \item ✅ Modern design (Tailwind CSS 4)
    \item ✅ Dark mode only (MVP)
    \item ⚠️ Keyboard navigation (basic; can improve)
    \item ✅ Responsive (desktop + tablet)
    \item ✅ Actionable error messages
    \item ⚠️ Real-time form validation (password strength; can expand)
    \item ✅ <5min onboarding for core features
\end{itemize}

\subsection{Maintainability}

\begin{itemize}
    \item ✅ ESLint + Prettier enforced
    \item ⏳ OpenAPI 3.1 spec (can add; API simple enough without)
    \item ⚠️ SQL migrations (manual; can automate)
    \item ⏳ Structured logging (Amplify CloudWatch available)
    \item ✅ README with setup + deployment
    \item ⏳ Dependency scanning (can enable)
\end{itemize}

%==============================================================================
\section{External Interfaces}
%==============================================================================

\subsection{Database Interface}

\begin{itemize}
    \item \textbf{Primary Storage:} PostgreSQL (auth, notes, metadata, vectors)
    \item \textbf{Secondary Storage:} AWS S3 (PDFs, Markdown files, exports)
    \item \textbf{Vector Search:} pgvector HNSW indexes
    \item \textbf{Full-Text Search:} PostgreSQL FULLTEXT (to implement)
\end{itemize}

Key tables:
\begin{itemize}
    \item \texttt{app.login} --- User accounts
    \item \texttt{app.notes} --- Note documents + search/embedding columns (extracted\_text, embedding, search\_vector, relative\_path)
    \item \texttt{app.tree\_items} --- Folder hierarchy (with sort metadata)
    \item \texttt{app.attachments} --- PDFs and files
    \item \texttt{app.pdf\_annotations} --- Highlights and notes on PDFs
    \item (Phase 2) \texttt{app.quiz\_questions}, \texttt{app.flashcards} --- Learning data
\end{itemize}

\textbf{Indexes (Phase 1):} GIN index on search\_vector (full-text), ivfflat on embedding (semantic), compound indexes on user\_id+created\_at and user\_id+updated\_at for sorting/filtering.

\subsection{API Endpoints}

\textbf{Authentication}
\begin{itemize}
    \item \texttt{POST /api/auth/register}, \texttt{/login}, \texttt{/logout}, \texttt{/forgot-password}, \texttt{/reset-password}
    \item All endpoints handle JWT tokens and session management
\end{itemize}

\textbf{Notes \& Files}
\begin{itemize}
    \item \texttt{GET /api/notes} --- List user's notes
    \item \texttt{POST /api/notes} --- Create note
    \item \texttt{GET /api/notes/:id}, \texttt{PATCH /api/notes/:id}, \texttt{DELETE /api/notes/:id}
    \item \texttt{POST /api/upload} --- Upload PDF or file
    \item \texttt{GET /api/tree?sort=alphabetical|recent} --- Get folder hierarchy with sorting
\end{itemize}

\textbf{Search (Phase 1)}
\begin{itemize}
    \item \texttt{GET /api/search?q=query\&type=fuzzy|semantic\&limit=20\&offset=0} --- Search by filename + content
    \item \texttt{POST /api/notes/:id/embed} --- Generate embedding for semantic search (internal/async)
    \item Returns: filename, relative path, excerpt, relevance score, creation/update dates
\end{itemize}

\textbf{RAG \& AI (Phase 2)}
\begin{itemize}
    \item \texttt{POST /api/rag/chat} --- Submit question, get answer with citations
    \item \texttt{GET /api/rag/search} --- Semantic search for RAG context
    \item \texttt{POST /api/notes/:id/index} --- Index PDF for RAG (background job)
\end{itemize}

\textit{See \texttt{SEARCH\_ARCHITECTURE\_PLAN.md} for full endpoint documentation with request/response examples.}

Full spec to come (or auto-generate from OpenAPI annotations).

\subsection{External Services}

\begin{itemize}
    \item \textbf{AWS S3:} File storage (PDFs, notes, exports)
    \item \textbf{AWS SES:} Transactional email (verification, reset)
    \item \textbf{AWS Amplify:} Deployment (auto-deploy from \texttt{prod} branch)
    \item \textbf{AWS RDS:} PostgreSQL database (managed backups)
    \item \textbf{OpenAI/Other LLMs:} User-supplied keys (not backend-managed)
\end{itemize}

%==============================================================================
\section{Data \& Compliance}
%==============================================================================

\subsection{Data Handling}

\begin{itemize}
    \item ✅ Encryption in transit (TLS 1.2+)
    \item ✅ Encryption at rest (S3 managed keys, RDS encryption)
    \item ⏳ 7-day soft-delete before hard-delete
    \item ✅ JWT tokens in \texttt{httpOnly} cookies
    \item ✅ Token expiry enforced (pwd reset: 1h; email verify: 24h)
    \item ✅ Vault export available (full data portability)
\end{itemize}

\subsection{Compliance}

\begin{itemize}
    \item ⚠️ \textbf{GDPR:} Basic compliance (data access, deletion, portability)
    \item ⚠️ \textbf{FERPA:} Not applicable unless LMS integration happens
    \item ✅ \textbf{MIT License:} Notea attribution in docs and UI
\end{itemize}

\subsection{Logging \& Audit}

\begin{itemize}
    \item ✅ API access logs (Amplify CloudWatch)
    \item ✅ Authentication events (login, reset, verify)
    \item ⏳ Document upload logging
    \item ⏳ Deletion logging with hashes (for recovery)
    \item ✅ 90+ day retention
\end{itemize}

%==============================================================================
\section{Verification \& Acceptance}
%==============================================================================

\subsection{Testing Strategy}

\begin{itemize}
    \item \textbf{Unit Tests:} Core logic (auth, SM-2, chunking); 70\%+ coverage
    \item \textbf{Integration Tests:} API endpoints with test DB
    \item \textbf{E2E Tests:} Critical flows (register → verify → upload → ask question)
    \item \textbf{Manual QA:} 3-5 test accounts, realistic usage
    \item \textbf{Performance:} Lighthouse audits (informative, not blocking)
    \item \textbf{Load Testing:} Vector search performance at scale (TBD)
\end{itemize}

\subsection{Acceptance Criteria}

The system is MVP-ready when:

\begin{enumerate}
    \item ✅ All ``Must Have'' auth features implemented
    \item ✅ Note creation, editing, sync working
    \item ✅ PDF upload and annotation working
    \item 🔄 RAG pipeline: chunking, embeddings, hybrid search, citations
    \item ⏳ Quiz generation (basic)
    \item ⏳ Flashcard system (basic SM-2)
    \item ✅ \texttt{docker-compose up} works fresh
    \item ✅ AWS Amplify auto-deploy functional
    \item ✅ Code coverage $\geq 70\%$ for core logic
    \item ✅ Academic stakeholder approval
\end{enumerate}

%==============================================================================
\section*{Appendix A: Project Phases}
%==============================================================================

\subsection*{Phase 1: Foundation + Search (Weeks 1-3)}

\textbf{Core Infrastructure (✅ Complete)}
\begin{itemize}
    \item ✅ Auth system (register, login, password reset)
    \item ✅ Note CRUD + folder structure
    \item ✅ PDF upload + annotations
    \item ✅ File sync to S3
    \item ✅ AWS Amplify deployment
    \item ✅ PostgreSQL + pgvector setup
\end{itemize}

\textbf{Search System (🔄 Planned/In Progress)}
\begin{itemize}
    \item 🔄 Database schema: new columns (extracted\_text, embedding, relative\_path, search\_vector)
    \item 🔄 Fuzzy full-text search endpoint (\texttt{GET /api/search?type=fuzzy})
    \item 🔄 Semantic search endpoint (\texttt{GET /api/search?type=semantic})
    \item 🔄 Embedding generation endpoint (\texttt{POST /api/notes/:id/embed})
    \item 🔄 Tree sorting (alphabetical, recent)
    \item 🔄 Search UI: Cmd+K overlay, relevance scoring, pagination
    \item ⏳ Phase 1.5: Date filtering UI (MVP 2)
\end{itemize}

\textit{See \texttt{SEARCH\_ARCHITECTURE\_PLAN.md} for full implementation roadmap and scalability considerations.}

\subsection*{Phase 2: RAG Pipeline (Weeks 3-5) --- 🔄 IN PROGRESS}
\begin{itemize}
    \item 🔄 PDF text chunking (500 tokens, 50-100 overlap)
    \item 🔄 OpenAI embedding integration
    \item 🔄 Vector similarity search for RAG context
    \item 🔄 RAG chat UI with citations
    \item 🔄 Streaming LLM responses with user-managed API keys
    \item 🔄 Citation formatting (page numbers, section headers)
\end{itemize}

\subsection*{Phase 3: Features \& Polish (Weeks 5-6) --- ⏳ PLANNED}
\begin{itemize}
    \item Quiz generation from notes/PDFs
    \item Flashcard system (SM-2 spaced repetition)
    \item LMS integration (Canvas OAuth + assignment sync)
    \item Calendar/timetable view
    \item Analytics dashboard (study time, weak topics)
    \item Manual QA testing with realistic scenarios
    \item Final demo + academic stakeholder approval
\end{itemize}

%==============================================================================
\section*{Appendix B: LLM Strategy (User-Managed Keys)}
%==============================================================================

\textbf{See separate document:} \texttt{docs/LLM\_STRATEGY.md}

Key points:
\begin{itemize}
    \item Users provide own API keys (stored in browser)
    \item Backend never stores or sees keys
    \item Supports OpenAI, Anthropic, LiteLLM, custom servers
    \item Team uses internal Qwen model during development (free)
    \item Frontend calls LLM directly; backend coordinates search only
\end{itemize}

%==============================================================================
\section*{Appendix C: Revision History}
%==============================================================================

\begin{longtable}{@{}llll@{}}
\toprule
\textbf{Version} & \textbf{Date} & \textbf{Change} & \textbf{Reason} \\
\midrule
\endhead
1.0 & Original & MariaDB, microservices, monolithic RAG & Initial specification \\
2.0 & 2025-03-06 & PostgreSQL + pgvector, monolith, Agile living doc & Team expertise; pragmatic MVP; rebrand to OghmaNotes \\
2.1 & 2025-03-06 & Comprehensive search architecture (Phase 1) & Planned fuzzy + semantic search with Cmd+K UI \\
\bottomrule
\end{longtable}

Future versions as tech stack evolves and phases complete.

%==============================================================================
\section*{Appendix D: Glossary}
%==============================================================================

\begin{itemize}
    \item \textbf{RAG:} Retrieval-Augmented Generation --- vector search + LLM context
    \item \textbf{pgvector:} PostgreSQL extension for vector embeddings
    \item \textbf{Embedding:} Numerical representation of text (1536-dim)
    \item \textbf{Chunking:} Splitting documents into segments for embedding
    \item \textbf{Hybrid Search:} Combining keyword (FULLTEXT) + semantic (vector) results
    \item \textbf{SM-2:} SuperMemo 2 algorithm for spaced repetition scheduling
    \item \textbf{PWA:} Progressive Web App --- installable web app with offline support
    \item \textbf{JWT:} JSON Web Token for stateless authentication
    \item \textbf{MVP:} Minimum Viable Product --- feature set for semester end
    \item \textbf{Agile:} Iterative development; adapt as reality evolves
\end{itemize}

%==============================================================================
\section*{Appendix E: About the Name}
%==============================================================================

\textbf{Oghma} (pronounced ``OH-muh'') is the Celtic deity of eloquence, learning, and writing from Irish mythology. A member of the Tuatha Dé Danann, Oghma created the Ogham script---the ancient Irish writing system---to preserve and transmit knowledge. He embodied both intellectual mastery and strength, serving as Champion while inventing systems for organizing language.

This dual nature of knowledge preservation and capability reflects our platform's mission: create structured learning systems that empower students through intelligent organization and retrieval of their study materials.

\vspace{2cm}

\noindent
\textit{Document Status: Living document --- updated as project evolves}  
\textit{Next Review: After RAG pipeline implementation (Week 4)}

\end{document}
